\documentclass[svenska]{style/LTHthesis}

% Preamble
% from template
\usepackage[T1]{fontenc}%
\usepackage{mathptmx, helvet}%
\usepackage{graphicx}%

\usepackage{wrapfig}
\usepackage{url}
\raggedbottom

\usepackage{xcolor}
\usepackage{hyperref}
\hypersetup{
  colorlinks,
  allcolors={blue!40!black},
}

\pagenumbering{gobble}
%
%%%%%%%%%%%%%%%%%%%%%%%%%
%%% Unified notation
%%%%%%%%%%%%%%%%%%%%%%%%%

%% Mathematics
% Shorthand
\newcommand{\R}{\mathbb{R}}
\newcommand{\T}{\mathrm{\small{T}}}

% Functions
\newcommand{\LL}{\mathcal{L}_2}
\newcommand{\funof}[1]{\left( #1 \right)}
\newcommand{\abs}[1]{\left|#1\right|}
\newcommand{\norm}[1]{\lVert#1\rVert}
\newcommand{\E}[1]{\mathbb{E}\left[ #1 \right]}
\newcommand{\trace}[1]{\mathrm{tr}\left(#1\right)}
\newcommand{\diag}[1]{\mathrm{diag}\left\{#1\right\}}

%% Control
\newcommand{\plant}{\mathcal{P}}
\newcommand{\ctrler}{\mathcal{C}}

%% Code

%% Real-Time
%

% Document
\begin{document}

\noindent
\begin{minipage}[]{0.25\textwidth}
    \noindent\includegraphics[width=20mm]{style/LUCsv}
    \par\vspace*{3mm} 
\end{minipage}\hfill
\begin{minipage}[]{0.74\textwidth}
    \vspace*{-3mm}
    \sffamily
    \begin{flushleft}
        \huge Vad är det Värsta som kan Hända när Datorn Krånglar? \\[2mm]
        \Large Nils Vreman\\[1mm]
        \normalsize Institutionen för Reglerteknik
    \end{flushleft}
\end{minipage}

\par\vspace{1mm}  

\begin{sffamily}\noindent
    Populärvetenskaplig sammanfattning av doktorsavhandlingen \emph{Analysis of Embedded Real-Time Control Systems Subject to Computational Overruns}, juni 2023.
    Avhandlingen kan laddas ner från: \url{http://www.control.lth.se/publications}
\end{sffamily}

\vspace{2mm}\hrule\vspace{3mm}

\noindent
Datorer i olika storlekar och styrka finns överallt runt omkring oss.
Allt från mobilen i din ficka till bilen du åker med till jobbet innehåller många små datorer, också kallade \emph{inbyggda system}.
Varje litet inbyggt system har uppgifter som den måste uträtta, t.ex., beräkna hastigheten på din bil.
I de flesta fall finns det krav och specifikationer på att beräkningen måste uträttas inom en viss tidsram.
Om bromsen trycks in förväntas bilen sakta ner direkt.
I verkligheten tar det lite tid för det inbyggda systemet att beräkna nedbromsningshastigheten innan kommandot kan skickas till bromssystemet.
Det är dock inget märkligt och systemet är oftast designat med den extra tiden det tar att uträtta hela uppgiften i åtanke.

\begin{wrapfigure}{r}{0.45\textwidth}
    \vspace*{-2mm}
    \begin{centering}
        \includegraphics[width=0.45\textwidth]{figs/microcontroller.jpg}
    \end{centering}\small\par\medskip
        \noindent\emph{Exempel på ett inbyggt system. På bilden ses en Raspberry Pi Pico där det svarta chippet i mitten (storlek $7\times 7$~mm) är en mikrokontroller, dvs., ett inbyggt system.}\footnotemark
    \vspace*{-2mm}
\end{wrapfigure}
\footnotetext{Laserlicht, CC BY-SA 4.0, \url{https://creativecommons.org/licenses/by-sa/4.0/}, via Wikimedia Commons}

Tyvärr så kan inbyggda system krångla, precis som en persondatorer kan.
Bland annat kan bakgrundsprogram, virus, och föråldrad hårdvara/programvara störa beräkningen.
Det kan försena beräkningen och i värsta fall medföra att uppgiften inte blir klar alls.
Vad händer då om det inbyggda systemet krånglar på ett sådant sätt att tidsramen för beräkningen inte respekteras?
Just den frågan besvaras i den här avhandlingen.

Många inbyggda system sitter i apparater där det inte gör så mycket ifall något krånglar.
Till exempel har mycket hemelektronik (så som moderna stereosystem, TV-apparater, och kylskåp) inbyggda system för att styra funktionaliteten och kommunicera med andra smarta system.
Det värsta som kan hända här är att (till exempel) videokvaliteten inte är helt perfekt på TVn.
Värre konsekvenser uppstår ifall det inbyggda systemet krånglar i ett bromssystem, en pacemaker, eller i ett flygplans styrsystem; då kan det uppstå verklig fara för människor i närheten och värstascenariot inkluderar även dödsfall.

Den här avhandlingen utvecklar matematiska verktyg och metoder för att förbättra analysen av säkerhetskritiska system, så som bromssystemet i en bil.
Eftersom det är livsviktigt att ingen person kommer till skada fokuseras analysen på att se till att det hela systemet är säkert, även under väldigt dåliga förhållanden för det inbyggda systemet.
Dåliga förhållanden innebär här, bland annat, att det inbyggda systemet inte lyckas beräkna något under en väldigt lång tid på grund av, t.ex., en hacker-attack.

\textbf{FORTSÄTT HÄR}

Att ett system är säkert i teorin innebär dock inte att det är det även i praktiken.


Här skall texten fortsätta, tills den blir antingen två eller (gärna) fyra sidor lång.
\end{document}
